\documentclass[12pt,a4paper]{article}
\usepackage[UTF8,heading=true]{ctex}
\usepackage{geometry}
\usepackage{amsmath,amssymb}
\usepackage{hyperref,fancyhdr,xcolor,setspace}
\usepackage{ragged2e}

\geometry{left=2.5cm,right=2.5cm,top=2.5cm,bottom=2.5cm}
\setlength{\headheight}{15pt}
\setstretch{1.5}
\setlength{\emergencystretch}{3em}
\hfuzz=20pt
\sloppy
\hypersetup{hidelinks}

\pagestyle{fancy}
\fancyhf{}
\fancyhead[C]{说明书摘要}
\fancyfoot[C]{\thepage}
\renewcommand{\headrulewidth}{0pt}

\begin{document}

\begin{center}
    {\Huge\heiti 中国移动专利申请} \\[0.5em]
    {\Huge\heiti 说明书摘要} \\[1em]
\end{center}

\noindent\textbf{摘要：}
本发明涉及硅基光子集成与量子密钥分发QKD技术领域，提出一种面向微型化终端的硅基光子QKD芯片封装与片上自校准方案。该方案在芯片上设置与量子主光路共模布设的参考波导/参考干涉结构，并配置温度与应力的片上传感单元以及加热器阵列、相位调制器等片上执行单元；在密钥分发时序中插入静默时隙或采用参考光功率门控，在不污染量子探测的前提下采集参考响应与传感信息，基于温度-应力联合漂移模型估计漂移状态并生成补偿控制量，闭环调节干涉平衡与耦合状态，实现对终端发热、姿态变化与封装残余应力引起漂移的在线抑制。由此可在体积与功耗受限场景下提升QKD稳定性与可工程化部署能力。

\vspace{0.6em}
\noindent\textbf{附图说明：}
图1为本发明一实施例的硅基光子QKD芯片封装结构与片上自校准闭环框图。

\vspace{0.2em}
\noindent\textbf{摘要附图：}
图1。

\end{document}

